\documentclass[11pt]{article}

\usepackage[margin=1in]{geometry}
\usepackage{amsmath, amssymb, amsthm}
\usepackage{algorithm, algorithmic}
\usepackage{hyperref}
\usepackage{booktabs}
\usepackage{natbib}

\title{CVaR / Leximin Reductions for Group-Fair Classification:\\
       A Light Plugin for \texttt{fairlearn}}
\author{hinanohart\\
        \texttt{https://github.com/hinanohart/cvar-leximin-fair}}
\date{2026}

\begin{document}
\maketitle

\begin{abstract}
We present \texttt{cvar-leximin-fair}, a small open-source library that adds
two group-fairness reductions to the \texttt{fairlearn} evaluation ecosystem.
\textsc{CVaRReduction} minimizes the Conditional Value-at-Risk of group losses
at quantile $\alpha$; as $\alpha \to 0$ it recovers the strict maximin objective.
\textsc{LeximinReduction} produces a lexicographic-minimum allocation across
groups by iteratively pegging the worst-off group's loss. Both reductions
re-use the cost-sensitive ERM strategy of~\citet{agarwal2018reductions} and
are scikit-learn--shaped so that they slot directly into the
\texttt{fairlearn.metrics.MetricFrame} evaluation flow. We position the
package as an algorithmic tool for Rawlsian / difference-principle aggregation
\emph{while explicitly acknowledging} Sen's capability-approach critique of
primary-goods framings as a constraint on interpretation.
\end{abstract}

\section{Introduction}
Distributional fairness in machine learning is increasingly evaluated under
group-level worst-case lenses, motivated both by foundational normative work
\citep{rawls1971theory, sen1979equality} and by regulatory developments
\citep{euaiact2024, nyc_ll144_2023}. Existing tooling around
\texttt{fairlearn}~\citep{bird2020fairlearn} and AIF360~\citep{bellamy2018aif360}
focuses on parity-style constraints (demographic parity, equalized odds)
rather than on \emph{worst-tail} or \emph{lexicographic} objectives.

This paper introduces \texttt{cvar-leximin-fair}, a Apache-2.0--licensed
plugin that fills this gap with a minimal, dependency-light implementation
of two reductions:
\begin{enumerate}
\item \textsc{CVaRReduction}: minimizes
$\mathrm{CVaR}_{\alpha}\!\left( \{ L_g \}_{g \in \mathcal{G}} \right)$
where $L_g$ is the per-group $0/1$ loss.
\item \textsc{LeximinReduction}: iteratively pegs each successive worst-off
group's loss, producing a leximin allocation in the limit.
\end{enumerate}

\section{Background}
\subsection{Reductions for fair classification}
\citet{agarwal2018reductions} cast group-fair classification as a
constrained minimization
$\min_{h \in \mathcal{H}} \mathbb{E}[\ell(h(X), Y)] \text{ s.t. } M h \le c$,
solved via Lagrangian reduction to cost-sensitive ERM. Their
\texttt{ExponentiatedGradient} and \texttt{GridSearch} routines underpin
\texttt{fairlearn.reductions}.

\subsection{CVaR and leximin}
The Conditional Value-at-Risk at level $\alpha$ of a discrete distribution
is the conditional mean of its upper $\alpha$-tail; it is a coherent risk
measure~\citep{rockafellar2000optimization}. The leximin order on loss
vectors is the lexicographic order on the sorted (descending) tuple, and
strictly refines the maximin order. Both are natural operationalizations
of the difference principle~\citep{rawls1971theory}.

\subsection{The Sen critique}
\citet{sen1979equality, sen2009idea} argues that Rawlsian primary-goods
aggregation is a form of \emph{resource fetishism}: equalizing resources
does not equalize \emph{capability}. We follow this critique: our package
equalizes \emph{loss}, which is a resource-proxy at best.
\textbf{Adopters should not interpret a CVaR- or leximin-fit model as
``capability-fair''.}

\section{Method}
\begin{algorithm}[h]
\caption{CVaR reduction (cost-sensitive reweighting)}
\begin{algorithmic}[1]
\STATE Initialize $w \gets \mathbf{1}$
\FOR{$t = 1, \dots, T$}
  \STATE $h_t \gets \mathrm{ERM}(X, y, w)$
  \STATE $L_g \gets \mathrm{loss}(h_t, X_g, y_g)$ for each group $g$
  \STATE $\mathcal{S}_t \gets$ top-$\lceil \alpha |\mathcal{G}| \rceil$ groups by $L_g$
  \STATE $w_i \gets (1 + \eta) w_i$ for $i$ with $s_i \in \mathcal{S}_t$;
         renormalize so $\sum_i w_i = n$
  \STATE Track $\mathrm{CVaR}_\alpha$; retain best-objective iterate
\ENDFOR
\STATE \textbf{return} best-iterate $h^\star$
\end{algorithmic}
\end{algorithm}

\noindent\textbf{Convergence note.} The iterative re-weighting is a
no-regret update with respect to the Lagrangian dual; on a fixed sample,
sub-linear convergence of the running average to a CVaR-optimal classifier
follows from~\citet{agarwal2018reductions} (Theorem 1) when
$\mathcal{H}$ is closed under convex combination. Best-iterate selection
gives \emph{weak} monotonicity vs.\ the initial unweighted fit (returned
$\mathrm{CVaR}_\alpha$ is no worse, but ties are possible and observed
on the Adult experiment).

\section{Experiments}
Adult Income (UCI / OpenML, $n = 48{,}842$, sensitive attribute = \texttt{sex}),
70/30 stratified split. Base estimator: \texttt{LogisticRegression(liblinear)}.

\begin{table}[h]
\centering
\begin{tabular}{lrrr}
\toprule
Method & Accuracy & Worst-group loss & $\mathrm{CVaR}_{0.5}$ \\
\midrule
Baseline (LR) & 0.8477 & 0.1903 & 0.1903 \\
CVaR ($\alpha=0.5$) & 0.8482 & 0.1899 & 0.1899 \\
Leximin & 0.8479 & 0.1902 & 0.1902 \\
\bottomrule
\end{tabular}
\caption{Test-set results on UCI Adult (sensitive attribute \texttt{sex},
70/30 stratified split, seed 0); values are those produced by the reference
notebook \texttt{examples/adult\_income.ipynb} at the v0.1.0 tag.
\textbf{CVaR improves slightly over baseline; leximin improves slightly
but less than CVaR}, which is expected on this benchmark where parity is
already nearly satisfied. Numbers update with re-execution.}
\end{table}

Adult \texttt{sex} is a known-easy fairness benchmark where parity is
already close to satisfied by the base estimator; the gains from CVaR /
leximin are small in this setting and the value is in the
\emph{audit-time} mechanism, not the train-time gain.

\section{Limitations and Ethics}
\begin{itemize}
\item Equalizing loss is not equalizing capability~\citep{sen1979equality}.
\item This package is not a compliance certification. Regulatory requirements
(EU AI Act Articles 9--17~\citep{euaiact2024}; NYC Local Law 144~\citep{nyc_ll144_2023})
define what is acceptable; our metrics may or may not satisfy them in any
specific jurisdiction at any given date.
\item Worst-case group definitions inherit the protected-attribute taxonomy
fed in; pre-processing of $s$ is out of scope.
\end{itemize}

\section{Conclusion}
We released a minimal Apache-2.0 plugin that brings CVaR and leximin
reductions to \texttt{fairlearn}. The implementation favors clarity over
optimization scale; pull requests with native \texttt{fairlearn.reductions.Moment}
subclasses are welcome.

\bibliographystyle{plainnat}
\bibliography{references}

\end{document}
